It is worth investigating whether there is a difference in results when considering only tumors or only metastases.
If the model performs substantially better on these subsets than on their union,
this is valuable on its own and provides a valuable new insight.
Precisely, it suggests that good rulesets for tumors contradict good rulesets for metastases.

Furthermore, a metastasis tends to have more genetic mutations than a tumor.
If more mutations cause more complicated and more unique behavior within the cell,
it may be harder to deduce general rules governing PD-L1 expression in metastases.
Again, focusing on one or the other independently could be useful in this case.

A related hypothesis is that metastases have higher PD-L1 expression due to their tendency for more mutations.
To investigate this, the distribution of PD-L1 expression per tissue type is shown in \autoref{fig:tissue_status_density_broad}.
Note that cell lines belonging to Papiloma or Normal are disregarded due to a lack of samples.
The density graphs have a similar shape.
The mean PD-L1 expression per tissue status is also similar.
Alternatively, if the hypothesis is true, we expect to see relatively many metastases among samples with high PD-L1 expression.
The tissue status distribution per PD-L1 expression quantile bin is depicted in \autoref{fig:tissue_status_quantile_broad}.
This histogram does not suggest the balance of tumors and metastases changes significantly with PD-L1 expression.
Thus, the results suggest that the hypothesis is false.

Still, there may be a difference in the behavior of the model between tumors and metastases.
To test this, the greedy method was used on the Broad dataset for tumors and metastases separately.
Both the tumors in Broad and the metastases in Broad were treated as separate datasets,
each undergoing independent discretization, and were subjected to the algorithm independently.
We consider $Q_y = 0.5, 0.6$ and $D=2$.
These values of $Q_y$ are the ones that were identified as most promising in \autoref{sec:boolean-results}.
The choice $D=2$ is practical: the greedy method is easier to use,
and for $D=2$, the greedy method is orders of magnitude faster than for higher $D$.
It is not purely practical;
in \autoref{sec:threshold-results} it was shown that performance barely increased when increasing $D$.
However, these results need not extend to this experiment.
It may be that other values of $Q_y$ and higher values of $D$ yield better results on these subsets.

The results are shown in \autoref{tab:greedy_type_tumor_yq_4_mw_0_C_3_D_2} through
\autoref{tab:greedy_type_metastasis_yq_5_mw_0_C_3_D_2}.
We note, without further analysis,
that first rules found in each of the four variants are all qualitatively different.

In all four cases, there is effectively no value in adding more than one rule.
Recall that after the first rule is found, all covered points are stripped away and the best rule for the remaining points
is found through enumeration.
Therefore, there is simply no effective rule for the leftover points.

Accuracy is slightly higher for tumors.
Therefore, the results of this experiment suggest that
PD-L1 expression is slightly more easily captured in a ruleset for tumors.
However, this says nothing about whether the rulesets are biologically meaningful.

Accuracy for metastases is lower than what is observed for the full dataset.
This suggests that contradicting rulesets between tumors and metastases are not the cause of low performance when using the full dataset.



\begin{figure}[h!]
    \centering
    \includesvg[width=0.8\textwidth]{gen/data/tissue_type_density_broad.svg}
    \caption{Density of PD-L1 expression for tissue types Tumor, Metastasis and Unknown in the Broad dataset.}
    \label{fig:tissue_status_density_broad}
\end{figure}


\begin{figure}[h!]
    \centering
    \includesvg[width=0.8\textwidth]{gen/data/tissue_type_quantile_broad.svg}
    \caption{Distribution of tissue type within various quantile bins of PD-L1 expression in the Broad dataset}
    \label{fig:tissue_status_quantile_broad}
\end{figure}

\clearpage

\vspace*{\fill}
\begin{table}[h]
    \centering
    \bwtable{gen/PDL1ex2/tab_greedy_type_tumor_yq_4_mw_0_C_3_D_2.tex}
    \caption{Results of the greedy method for tumors, $Q_y = 0.5$.}
    \label{tab:greedy_type_tumor_yq_4_mw_0_C_3_D_2}
\end{table}

\begin{table}[h]
    \centering
    \bwtable{gen/PDL1ex2/tab_greedy_type_tumor_yq_5_mw_0_C_3_D_2.tex}
    \caption{Results of the greedy method for tumors, $Q_y = 0.6$.}
    \label{tab:greedy_type_tumor_yq_5_mw_0_C_3_D_2}
\end{table}
\vspace*{\fill}

\clearpage

\vspace*{\fill}
\begin{table}[h]
    \centering
    \bwtable{gen/PDL1ex2/tab_greedy_type_metastasis_yq_4_mw_0_C_3_D_2.tex}
    \caption{Results of the greedy method for metastases, $Q_y = 0.5$.}
    \label{tab:greedy_type_metastasis_yq_4_mw_0_C_3_D_2}
\end{table}

\begin{table}[h]
    \centering
    \bwtable{gen/PDL1ex2/tab_greedy_type_metastasis_yq_5_mw_0_C_3_D_2.tex}
    \caption{Results of the greedy method for metastases, $Q_y = 0.6$.}
    \label{tab:greedy_type_metastasis_yq_5_mw_0_C_3_D_2}
\end{table}
\vspace*{\fill}